% Avsnitt 11.7, ur lectures/L11/README.md.

\section{Sammanfattning}
\label{c11:sec:review}

\needspace{6\baselineskip}
\subsection*{Det här ska du kunna nu}
\begin{itemize}
  \item Rita kontrollerns blockschema och säga vad varje block ansvarar för, och lika viktigt, vad
    det inte ansvarar för.
  \item Skriva toppnivåns entitet med rätt portordning och rätt typer.
  \item Förklara varför arkitekturen är tom nu och vilket kapitel som fyller vilken del av den.
  \item Köra GHDL:s två första steg på den tomma toppnivån, läsa vad vart och ett rapporterar, och
    säga vad det tredje steget kräver som ännu inte finns.
  \item Säga vad positionell portbindning köper och vad den ger upp.
\end{itemize}

\needspace{6\baselineskip}
\subsection*{Frågor att testa dig själv med}
\begin{enumerate}
  \item Varför äger \code{bit_timer} bara \emph{när}, och inte \emph{vad}?
  \item Vilket block känner till CAN-ramens format, och varför bara ett?
  \item Varför synkroniserar \code{can_controller} sina egna asynkrona ingångar i stället för att
    kräva att den som instansierar den gör det?
  \item Vad går sönder om du byter plats på två portar av samma typ, och när skulle du märka det?
  \item Vad gör \code{ghdl -a}, \code{ghdl -e} och \code{ghdl -r}, var för sig?
  \item Vilken registerkarterad operation svarar mot \code{tx_req}, och vilken mot \code{tx_done}?
\end{enumerate}

\needspace{6\baselineskip}
\subsection*{Riktvärde}
Gruppens repo bör finnas, med skyddad \code{main}, och de utdelade filerna bör ha nått \code{main}
via en granskad Pull Request. Toppnivåns entitet bör finnas och analysera rent. Därefter är det upp
till gruppen: någon kan mycket väl börja på \code{bit_timer} eller \code{crc15} redan nu. De fyra
lövmodulerna, \code{bit_timer}, \code{crc15}, \code{tx_shift_reg} och \code{rx_shift_reg}, beror
inte på varandra och har var sin egen testbänk, så en grupp om 4--5 kan bygga dem parallellt. Det
är den enda punkt i projektet där ordningen på riktigt spelar roll för hur snabbt ni går framåt; se
bilaga~\ref{app:project}.

\needspace{6\baselineskip}
\subsection*{Blicken framåt}
\Chapref{c12:ch} bygger de två första byggblocken som går in i den tomma arkitekturen, båda små:
synkroniseraren \code{meta_prev}, som inte vet någonting alls, och \code{bit_timer}, det första
blocket som vet något om CAN.
